% ભાગ: first-order-logic
% પ્રકરણ: beyond
% વિભાગ: introduction

\documentclass[../../../include/open-logic-section]{subfiles}

\begin{document}

\olfileid{fol}{byd}{int}

\olsection{વિહંગાવલોકન}

પ્રથમ-ક્રમ તર્કશાસ્ત્ર જ રસપ્રદ એકમાત્ર તાર્કિક તંત્ર નથી: તેના અનેક
વિસ્તારો અને રૂપાંતરો છે. સામાન્ય રીતે કોઈ તાર્કિક તંત્રમાં ભાષાનું
વિધિવત્ નિર્દેશન, મોટેભાગે---પણ હંમેશાં નહિ---નિષ્પત્તિ-તંત્ર, અને
મોટેભાગે---પણ હંમેશાં નહિ---અભિપ્રેત અર્થવિજ્ઞાન હોય છે. પરંતુ આ
શબ્દનો તકનિકી ઉપયોગ એક સ્વાભાવિક પ્રશ્ન ઊભો કરે છે: પ્રથમ-ક્રમ
તર્કશાસ્ત્રથી જુદાં તાર્કિક તંત્રોનો સહજ અથવા તત્ત્વચિંતનાત્મક અર્થમાં
વપરાતા ``તર્ક'' શબ્દ સાથે શું સંબંધ છે? નીચે વર્ણવેલાં બધાં તંત્રો કોઈ
ને કોઈ પ્રકારના તર્કવિચારનું નિદર્શન કરવા માટે રચાયાં છે; તેમને
તાર્કિક શું બનાવે છે તે આપણે કહી શકીએ?

તેનો કોઈ સહેલો જવાબ મળતો નથી. ``તર્ક'' શબ્દ જુદી રીતે અને જુદા
સંદર્ભોમાં વપરાય છે, અને ``સત્ય''ના ખ્યાલની જેમ આ ખ્યાલનું પણ અનેક
તત્ત્વચિંતનાત્મક દૃષ્ટિકોણોથી વિશ્લેષણ થયું છે. દાખલા તરીકે, તાર્કિક
તર્કવિચારનો ધ્યેય કયાં વિધાનો અનિવાર્યપણે સત્ય છે, પૂર્વાનુભવી રીતે
સત્ય છે, અતાર્કિક પદોના અર્થઘટનથી સ્વતંત્રપણે સત્ય છે, તેમના સ્વરૂપને
કારણે સત્ય છે, અથવા ભાષાકીય રૂઢિથી સત્ય છે તે નક્કી કરવાનો છે એવું
કોઈ માની શકે; અને આ દરેક કલ્પનાને ઘણું સ્પષ્ટીકરણ જોઈએ. ધ્યાન માત્ર
ગણિતમાં વપરાતા તર્કશાસ્ત્ર પૂરતું મર્યાદિત કરીએ તોપણ તેની વ્યાપ્તિ
વિશે બહુ ઓછી સંમતિ છે. દાખલા તરીકે, \textit{Principia Mathematica}માં
રસેલ અને વ્હાઇટહેડે ફ્રેગેથી શરૂ થયેલી {\em તર્કવાદી} પરંપરામાં
તર્કશાસ્ત્રના પાયા પર ગણિત વિકસાવવાનો પ્રયત્ન કર્યો હતો. તેમનું
તાર્કિક તંત્ર નીચે વર્ણવેલા તંત્ર જેવું ઉચ્ચ-પ્રકાર તર્કશાસ્ત્રનું એક
સ્વરૂપ હતું. આખરે તેમને એવી સ્વયંસિદ્ધિઓ દાખલ કરવાની ફરજ પડી જે મોટા
ભાગનાં ધોરણો મુજબ શુદ્ધ તાર્કિક લાગતી નથી (ખાસ કરીને અનંતતાનું
સ્વયંસિદ્ધિ અને લઘુનીયતાનું સ્વયંસિદ્ધિ); તેમ છતાં ઉચ્ચ-ક્રમ તર્કવિચારનાં
કેટલાંક સ્વરૂપોને તાર્કિક ગણવાં જોઈએ એવું કોઈ માની શકે. તેનાથી ઊલટું,
ક્વાઇનનું તત્ત્વમીમાંસક માળખું ``વિધાનો''ને ચર્ચાની કાયદેસર વસ્તુઓ
ગણતું નથી; તેથી તેઓ દલીલ કરે છે કે દ્વિતીય-ક્રમ અને ઉચ્ચ-ક્રમ
તર્કશાસ્ત્ર ખરેખર તર્કશાસ્ત્રના વેશમાં આવેલાં ગણસિદ્ધાંતનાં સ્વરૂપો છે.
બીજા શબ્દોમાં, વિધેયો પર પરિમાણીકરણ ધરાવતાં તંત્રો શુદ્ધ તાર્કિક નથી.

હાલ પૂરતું આવા તત્ત્વચિંતનાત્મક પ્રશ્નોને બીજા દિવસ માટે મૂકી રાખવા
અને નીચેનાં તંત્રોને તાર્કિક હોય કે ન હોય એવા વિવિધ પ્રકારના
તર્કવિચારનાં વિધિવત્ આદર્શીકરણો તરીકે વિચારવું શ્રેષ્ઠ છે.

\end{document}
